%-----------------------------------------------------------------------
\section{Related Work}
%-----------------------------------------------------------------------

\paragraph{Badminton datasets.}  ShuttleSet provides broadcast videos, set-level
annotations, stroke labels, homography information, and pre-extracted features
for badminton analysis \cite{shuttleset}. Its paired structure—raw video
alongside structured labels—makes it well suited to evaluating how much tactical
information survives end-to-end prediction from video. BadmintonDB focuses on
player-specific match analysis and outcome prediction \cite{badmintondb}.
Recently, the FineBadminton project introduced a large-scale dataset with a
multi-level annotation hierarchy ranging from low-level actions to decision
evaluation \cite{finebadminton}. The dataset and associated FBBench benchmark
are designed to challenge multimodal models on nuanced spatio-temporal reasoning,
but their focus is on higher-level semantics rather than fine-grained stroke
classification. We adopt ShuttleSet as our primary benchmark because its
annotations support evaluation at the stroke, rally, and sequence levels
simultaneously.

\paragraph{Rally filtering, shuttle tracking, and hit detection.}
Rally filtering separates playable footage from non-rally broadcast segments, a
prerequisite for any downstream processing. Shuttle tracking is typically
formulated as frame-level heatmap or coordinate prediction; TrackNetV2
introduced an efficient network architecture for this task in badminton video
\cite{tracknetv2}, and we build on a TrackNetV3 implementation for coordinate
extraction and denoising. Hitting-event detection has been studied directly
using specialized models \cite{chen2023hitting}. Recent work has also proposed
refining shuttle hit times by combining TrackNet trajectory predictions with
swing detection from YOLOv7 \cite{hsu2024enhancing}, improving shot extraction
in monocular-camera settings. In the present work, exact ShuttleSet hit-frame
labels define the unseen-video benchmark; trajectory-based hit detection
remains the prerequisite step for fully unlabeled deployment.

\paragraph{Pose estimation and stroke classification.}
Player pose provides body keypoints and on-court position, both of which are
necessary inputs to stroke classification. MMPose is a widely used
pose-estimation toolbox \cite{mmpose}; our pipeline employs a YOLO-pose
checkpoint for efficient per-window feature extraction. For stroke recognition,
the Badminton-Stroke Transformer (BST) models racket-sport strokes jointly from
player skeletons, positions, and shuttle kinematics \cite{chang2025bst}.
We fine-tune the \texttt{BST\_CG\_AP} variant as the final stroke classifier.

\paragraph{Sports tactic mining and immersive analysis.}
Tactical analysis from structured event sequences is an established area in
sports analytics. Stroke distributions characterize match composition,
transition matrices encode sequential stroke tendencies, and frequent-sequence
mining identifies recurring tactical phrases \cite{shuttleset}. Movement
profiling connects spatial displacement with stroke selection. KMeans clustering
\cite{kmeans} paired with silhouette analysis \cite{silhouette} offers a
standard approach for profile summarization, but cluster validity must be
confirmed through stability tests rather than visual inspection alone. Beyond
batch analysis, immersive systems such as VIRD provide interactive 3D visual
analytics for coaches and players by reconstructing rally scenes in virtual
reality \cite{vird2023}. Our work differs by focusing on automated recovery of
stroke sequences from raw video and on producing interpretable mining outputs
from predicted event streams rather than on visualization tools.

\paragraph{Positioning \method{} against prior work.}
Table~\ref{tab:comparison} summarises how \method{} sits relative to the
systems above across five dimensions: input modality, output, degree of
automation, support for tactical mining, and evaluation protocol. No prior
system simultaneously takes raw video as input, produces a structured stroke
event stream, \emph{and} validates the downstream mining against ground truth.
TrackNet-based systems handle the video-to-trajectory step but stop at shuttle
coordinates. BST handles stroke classification but requires pre-extracted
features, not raw frames. ShuttleSet's own analysis tools operate on
ground-truth annotations rather than predicted events, providing an upper bound
but not an end-to-end pipeline. VIRD delivers immersive visualisation but
relies on structured input and does not address recovery from raw video.
\method{} fills the gap between raw video and a validated tactical analysis by
chaining these components and measuring how much tactical signal survives
end-to-end prediction.

\begin{table*}[t]
  \caption{Qualitative comparison of related badminton analytics systems.
  ``Partial'' automation means the system processes raw video for some stages
  but relies on supplied metadata (hit-frame indices or homographies) for
  others.}
  \label{tab:comparison}
  \centering
  \small
  \setlength{\tabcolsep}{4pt}
  \renewcommand{\arraystretch}{1.15}
  \begin{tabular}{p{2.8cm} p{3.2cm} p{2.8cm} p{1.6cm} p{1.6cm} p{3.5cm}}
    \toprule
    System & Input & Output & Fully auto? & Tactical mining? & Evaluation \\
    \midrule
    ShuttleSet analysis \cite{shuttleset}
      & Ground-truth stroke annotations
      & Tactical statistics
      & No
      & Yes
      & Annotation-derived; no pred-vs-truth gap \\[3pt]

    BST \cite{chang2025bst}
      & Pre-extracted features (pose, shuttle)
      & Stroke class per window
      & No
      & No
      & Classification accuracy on held-out clips \\[3pt]

    TrackNet-based systems \cite{tracknetv2,hsu2024enhancing}
      & Raw broadcast video
      & Shuttle trajectory; hit-frame cues
      & Partial
      & No
      & Tracking coverage; hit-detection F1 \\[3pt]

    VIRD \cite{vird2023}
      & Structured / visual input
      & Interactive 3-D VR visualisation
      & No
      & Visualisation only
      & User study / visual inspection \\[3pt]

    \textbf{\method{} (ours)}
      & Raw video + supplied hit indices \& homographies
      & Stroke event stream + five mining analyses
      & Partial
      & Yes
      & Pred-vs-truth mining validation (KL, Frobenius, Spearman, ARI) \\
    \bottomrule
  \end{tabular}
\end{table*}